% Formalización reunida el 26 de septiembre de 2026.
% Procedencia y condiciones: COMUN/procedencia_20260926/.
\clearpage
\section{Normalisation conjointe de l’action, de la longueur et du couplage gravitationnel}
\label{md:rigidez}

La détermination d’une grandeur est conservée lors de la composition des lectures du
même état. On réunit ici le retour inscrit, la section d’action, la
vitesse constitutive et le caractère positif de route. L’antécédent commun
est la génération APP--TRIT--TPK, avec ses feuillets, son orientation, sa retenue, sa mémoire
et ses prolongements ; les coordonnées régionales et $\alpha$ interviennent comme
sorties de cette génération. La construction suivante utilise ces sorties
avant d’effectuer la reconnaissance dimensionnelle de $G$.

Le problème de normalisation comprend trois opérations distinctes : produire
une longueur à partir du retour enregistré, transporter le caractère sans perdre
son complément et comparer les lectures d’énergie et de longueur qui en résultent.
Les identités de cette section fixent conjointement ces opérations. Leurs
preuves portent sur des fibres finies ou sur des opérateurs bornés uniformément
positifs ; le secteur nul conserve son traitement séparé.

\subsection{Retour inscrit et section longitudinale}
\label{md:rigidez:longitud}

Dans le support marqué du transport, on conserve le support
\[
 \Omega=\{(j,k,q,u,v):j\in\mathbb Z/12\mathbb Z,\ 
 k-j\in\{0,3,6,9\},\ 1\le q\le8,\
 u,v\in\{1,2,4,5,7,8\},\ u<v\}.
\]
L’indice $j$ conserve la phase dodécaphasique ; les quatre décalages
conservent le quart de tour et les indices restants l’incidence marquée.
Par conséquent, $N=|\Omega|=12\cdot4\cdot8\binom62=5760$ compte les positions
d’incidence. L’espace de ce registre est transporté avec la mémoire
généalogique ; sa dimension finie correspond à ce support particulier.

Soient $B:V_U\to V_K$ l’inverseur de Hadamard transporté,
$B^*B=BB^*=I/4$, $u_K=N^{-1/2}(1,\ldots,1)$,
$P_K=u_Ku_K^*$ et $Q_K=I-P_K$. Nous définissons
\[
 C_0=Q_KB,\qquad M_0=P_KB,\qquad U=2B.
\]
L’orthogonalité des deux projecteurs démontre
\[
 C_0^*C_0+M_0^*M_0=I/4,\qquad C_0^*M_0=0,\qquad U^*U=UU^*=I.
\]
Ainsi, $x\mapsto(2C_0x,2M_0x)$ conserve l’archive complète de ce
transport. $U$ est sa normalisation unitaire ; l’évolution enrichie
$U_t$ maintient son propre domaine et sa mise à jour.

Soit $E_i=e_ie_i^*$ la résolution du registre et
$\Delta_\Omega(A)=\sum_iE_iAE_i$ son espérance diagonale.

\begin{lemma}[Retour longitudinal inscrit]
\label{md:rigidez:retorno}
Le retour et sa lecture inscrite satisfont
\[
 C_0C_0^*=Q_K/4,\qquad
 \Delta_\Omega(C_0C_0^*)=r_NI,\qquad
 r_N=\frac{N-1}{4N}=\frac{5759}{23040}.
\]
Avec la norme scalaire locale $n_H=\alpha^{16}$ du conoyau de $H_4$,
la composition longitudinale sur une cochaîne de longueur $\beta_*$ est
\[
 \mathscr T_L(\beta_*)=
 n_H\bigl(\Delta_\Omega(C_0C_0^*)\otimes I\bigr)\beta_*
 =q_\alpha\beta_*,
 \qquad q_\alpha=\alpha^{16}r_N.
 \label{md:rigidez:qalpha}
\]
\end{lemma}
\begin{proof}
L’égalité $BB^*=I/4$ donne $C_0C_0^*=Q_K/4$. Chaque entrée diagonale
de $P_K$ vaut $1/N$, ce qui prouve l’espérance. Multiplier une
cochaîne par ce scalaire conserve son type longitudinal et donne la dernière
égalité. Si $\beta_*(\bar a)=-U_a\beta_*(a)$, l’image satisfait la
même inversion. De même, à partir de
$\beta_m(a)=\sum_jU_j^{-1}\beta_{m+1}(a_j)$, on obtient l’identité
raffinée en multipliant les deux membres par $q_\alpha$.
\end{proof}

L’espérance possède une réalisation inscrite explicite :
$We_i=e_i\otimes e_i$ et
$W^*(A\otimes I)W=\Delta_\Omega(A)$. Pour
$R_K=Q_K/4$ et $P_W=WW^*$, le complément
$Z_\perp=(I-P_W)(R_K\otimes I)W$ vérifie
\[
 (R_K\otimes I)W=Wr_N+Z_\perp,\qquad W^*Z_\perp=0,\qquad
 Z_\perp^*Z_\perp=\frac{N-1}{16N^2}I.
\]
En effet, $R_K^2=R_K/4$, de sorte que le carré de la norme complète
est $r_N/4$, tandis que la composante inscrite a pour carré $r_N^2$.
La différence est précisément le complément écrit. La conservation de
ce complément accompagne la lecture longitudinale ; l’ajouter à une autre lecture
changerait l’observable.

L’espérance est unique parmi les applications bimodulaires vers
l’algèbre diagonale qui fixent ses projecteurs : sur une unité matricielle
$E_{ij}$, la bimodularité exige
$\mathcal E(E_{ij})=E_i\mathcal E(E_{ij})E_j$, qui vaut zéro si $i\ne j$,
et $\mathcal E(E_{ii})=E_{ii}$. Cette unicité conserve la résolution marquée.

Avec la section longitudinale de référence $L_*$ déjà spécifiée, il reste
\begin{equation}
 L_\alpha=q_\alpha L_*,\qquad D=L_\alpha U.
 \label{md:rigidez:longitud-def}
\end{equation}
Le lecteur utilise le retour comme opérateur linéaire sur une longueur. Une
lecture quadratique de la source produirait un autre type de grandeur. La
section $L_*$, la résolution inscrite et la norme locale sont des antécédents
explicites de cette composition.

\subsection{Transport polaire et réciprocité du même caractère}
\label{md:rigidez:polar}

Soit $X$ le caractère opératoriel du secteur positif non nul, en conservant la route et
ses normalisations. En dimension finie, on exige $X>0$ ; en dimension
infinie, on exige $X$ borné, autoadjoint et $X\ge\epsilon I$ pour un certain
$\epsilon>0$. Nous écrivons $Y=UXU^*$. La représentation radiale polaire est
\[
 R=(DX^2D^*)^{1/2}.
\]
La racine positive est unique et, en substituant $D=L_\alpha U$, on obtient
$R=L_\alpha Y$. De manière équivalente,
$\|Rv\|^2=\|XD^*v\|^2$ pour tout $v$ : la polarisation retrouve
$R^2=DX^2D^*$ et la positivité détermine $R$. Cette condition spécifie
la métrique induite qui est conservée.

L’action $\hbar$ et la vitesse constitutive $c$ appartiennent à la même
section. La réalisation circulaire compatible fixe
\begin{equation}
 \omega_0=\frac{c}{L_\alpha},\qquad
 t_0=\frac{\pi_{\rm HMT}L_\alpha}{54c},\qquad
 R_{108}=\frac{54}{\pi_{\rm HMT}}ct_0=L_\alpha.
 \label{md:rigidez:reloj}
\end{equation}
La longueur a été construite avant d’exprimer l’horloge circulaire. Cette
égalité permet de conserver la notation $R_{108}$ dans les développements
ultérieurs et fixe sa normalisation. Les bases dimensionnelles sont transportées
conjointement ; $L_*$ et la base d’arête $u_L$ ont des fonctions différentes.

\begin{theorem}[Détermination conjointe du couplage radial]
\label{md:rigidez:teorema}
Les lecteurs sur le même caractère,
\[
 H=\frac{\hbar c}{L_\alpha}Y,\qquad
 \Lambda=L_\alpha Y^{-1},\qquad R=L_\alpha Y,\qquad M=H/c^2,
\]
satisfont
\begin{equation}
 H\Lambda=\hbar cI,\qquad R\Lambda=L_\alpha^2I,\qquad
 R=\frac{L_\alpha^2}{\hbar c}H.
 \label{md:rigidez:identidades}
\end{equation}
Dans l’identification dimensionnelle de $R$ comme rayon gravitationnel réduit,
$R=GM/c^2$, le coefficient commun est unique :
\begin{equation}
 \boxed{G=\frac{c^3L_\alpha^2}{\hbar}.}
 \label{md:rigidez:G}
\end{equation}
\end{theorem}
\begin{proof}
La multiplication des lecteurs donne les deux premières identités.
Les multiplier à droite par $\Lambda^{-1}$ prouve la troisième.
Substituer $M=H/c^2$ identifie le coefficient. Si $G'$ conserve les mêmes
lectures, $(G'-G)M/c^2=0$ ; dans la fibre non nulle, l’inversibilité de $M$
implique $G'=G$.
Tous les produits sont définis dans la fibre complète sous les
hypothèses de l’énoncé.
\end{proof}

La preuve s’applique à tous les modes positifs de cette réalisation, sans choisir
la masse électronique pour normaliser $G$. L’identification du rayon
réduit est explicite ; le rayon $2GM/c^2$ correspond à une autre convention
de lecture de la même constante. Le secteur nul est maintenu en dehors de
l’inverse. L’extension aux opérateurs non bornés requiert un développement de
domaines et de limites distinct du présent théorème.

Le carré de $L_\alpha$ procède du produit des longueurs conjuguées.
L’action apparaît inversement lorsqu’on élimine $\Lambda$ entre l’énergie et la
longueur. Le passage de l’énergie à la masse et de la masse au rayon apporte $c^4$,
qui, combiné avec le $c$ de $H\Lambda=\hbar cI$, produit $c^3$.
Les deux normalisations fixent le facteur adimensionnel à un. L’analyse
dimensionnelle vérifie les exposants ; la composition détermine le coefficient.

\subsection{Échelles de Planck et spécialisation du caractère unitaire}
\label{md:rigidez:escalas}

La longueur construite, la section d’action et la vitesse déterminent
\[
 t_P=\frac{L_\alpha}{c},\qquad
 m_{P,a}=\frac{\hbar_a}{cL_\alpha},\qquad
 E_{P,a}=\frac{\hbar_a c}{L_\alpha}.
\]
Substituer \eqref{md:rigidez:G} dans
$\sqrt{\hbar_aG_a/c^5}$, $\sqrt{\hbar_ac/G_a}$ et
$\sqrt{\hbar_ac^5/G_a}$ démontre respectivement ces identités ;
la positivité fixe chaque racine. La longueur précède la reconnaissance
conventionnelle de ces échelles.
Pour un mode de caractère $y>0$,
\[
 m(y)=m_{P,a}y,\qquad r_g(y)=L_\alpha y,\qquad
 \bar\lambda(y)=L_\alpha/y.
\]
L’égalité des deux longueurs équivaut à $y=y^{-1}$ et, par conséquent,
à $y=1$. Le produit surfacique est conservé pour tous les modes positifs ;
la coïncidence des rayons caractérise le mode unitaire.

Le pas chronologique satisfait $t_0=(\pi_{\rm HMT}/54)t_P$ ;
la période de phase est $108t_0=2\pi_{\rm HMT}t_P$ et
$\omega_{108}=1/t_P$. Un horizon $R_h=\zeta_hL_\alpha$ conserve
sa condition géométrique propre. Pour la spécialisation de Schwarzschild
de la même masse, $R_h=2r_g=2L_\alpha y$.
En comparant des sections avec $L_\alpha,c$ fixes et
$\rho_a=\hbar_b/\hbar_a$, on obtient
$G_b=G_a/\rho_a$, $m_{P,b}=\rho_a m_{P,a}$ et
$E_{P,b}=\rho_a E_{P,a}$. Cette comparaison de sections conserve
$L_\alpha$ et $t_P$.

\subsection{Aire et restitution des secteurs de mémoire}
\label{md:rigidez:area}

L’incidence de bord conserve le regroupement de six trits.
Si $x_j=r_{2j-1}+3r_{2j}$, $r_i\in\{0,1,2\}$, l’application
\[
 (x_1,x_2,x_3)\longmapsto
 (r_1+3r_2+9r_3,\ r_4+3r_5+9r_6)
\]
est une bijection entre $(\mathbb Z/9\mathbb Z)^3$ et
$(\mathbb Z/27\mathbb Z)^2$. Son inverse extrait les six chiffres dans l’ordre,
en conservant l’endroit où la retenue est enregistrée. Dans le tore de bord,
le bloc $A=\{0,\ldots,8\}^2$ a deux collections de liens :
\begin{align*}
 \partial A_{90}&=\{((-1,j),(0,j)),((8,j),(9,j)):0\le j\le8\},\\
 \partial A_{120}&=\{((i,-1),(i,0)),((i,8),(i,9)):0\le i\le8\}.
\end{align*}
Chaque collection contient dix-huit liens. Leurs étiquettes conservent la
provenance de canal ; chaque lien porte un opérateur nombre de spectre
$\{0,1,2\}$. Soient $N_{90},N_{120}$ les sommes correspondantes.

L’échelle angulaire surfacique conserve le contre-angle déjà construit :
\[
 \theta_{\rm clk}=\frac{\pi C^*_{\deg}}{1080},\qquad
 a_0=\frac{1+2\cos(\pi/18)}3\,\theta_{\rm clk}.
\]
Le préfacteur est la trace réelle normalisée de
$\operatorname{diag}(1,e^{i\pi/18},e^{-i\pi/18})$.
Ainsi, $a_0$ conserve la signification de normalisation adimensionnelle d’aire.
La lecture d’aire est prise sur la branche positive
$\theta_{\rm clk}>0$ et $\gamma>0$.
Avec la fonctionnelle $\gamma$ de Barbero--Immirzi et la longueur déjà déterminée,
\[
 \mathcal A=\gamma a_0L_\alpha^2(N_{90}+N_{120}).
\]
Les trente-six opérateurs locaux agissent dans des facteurs tensoriels
distincts. Par conséquent, le spectre de $N=N_{90}+N_{120}$ est
$\{0,\ldots,72\}$ et la multiplicité de $n$ est le coefficient de $z^n$
dans $(1+z+z^2)^{36}$. La preuve consiste à multiplier les fonctions
génératrices des facteurs. La nouvelle longueur transporte l’unité
surfacique et conserve cette combinatoire.

\begin{proposition}[Lecture conjointe de l’aire et de la provenance sectorielle]
Le couple
$N=N_{90}+N_{120}$ et $\mathcal D_{\partial}=90N_{90}+120N_{120}$
restitue les deux secteurs :
\[
 N_{90}=4N-\frac{\mathcal D_{\partial}}{30},\qquad
 N_{120}=\frac{\mathcal D_{\partial}}{30}-3N.
\]
\end{proposition}
\begin{proof}
La matrice de la transformation est
$\left(\begin{smallmatrix}1&1\\90&120\end{smallmatrix}\right)$,
de déterminant $30$. Son inverse donne les formules. La commutativité
des deux opérateurs nombre implique celle des deux lecteurs.
\end{proof}
Dans la réalisation modulaire avec
$\rho_A=Z_A^{-1}\exp(-\Delta_{\rm mod}\mathcal D_{\partial})$,
$\Delta_{\rm mod}>0$, le calcul fonctionnel donne
$-\log\rho_A=(\log Z_A)I+\Delta_{\rm mod}\mathcal D_{\partial}$.
Les normalisations étant conservées, le couple formé par l’aire et le
Hamiltonien modulaire retrouve les deux occupations. La preuve de restitution
vaut pour chaque valeur construite de $\Delta_{\rm mod}$ ; le nouveau lecteur
longitudinal ne redéfinit ni ce coefficient ni l’état modulaire.
La somme des occupations et la provenance sectorielle expriment des observables
distincts. De même, les trente-six liens de réponse et les
quatre-vingt-une capacités unitaires de la coupure cubique conservent leurs fonctions.

\subsection{Variation conjointe et réduction compatible de mémoire}
\label{md:rigidez:variaciones}

Pour des collections différentiables en norme et uniformément positives dans un
voisinage, soit $\kappa=L_\alpha^2/(\hbar c)$. Alors
\[
 \delta R=\delta\kappa\,H+\kappa\,\delta H,\qquad
 \frac{\delta G}{G}=3\frac{\delta c}{c}
 +2\frac{\delta L_\alpha}{L_\alpha}-\frac{\delta\hbar}{\hbar}.
\]
La première identité résulte de la différentiation de
$R=L_\alpha Y$ et $H=(\hbar c/L_\alpha)Y$ ; elle ne requiert pas
$[Y,\delta Y]=0$. Une variation admissible de connexion qui agit dans $Y$
et conserve les trois sections maintient $\delta G=0$.
L’affirmation comprend les variations de cette réalisation radiale ;
l’équivalence entre des fonctionnelles de champ complètes a ses propres
domaines et conditions aux limites.

Soit maintenant
$Y=\left(\begin{smallmatrix}A&C\\C^*&D_i\end{smallmatrix}\right)>0$
dans la séparation entre lecture et mémoire, avec
$S=A-CD_i^{-1}C^*>0$. L’extension minimisante d’un vecteur $b$ est
$(b,-D_i^{-1}C^*b)$, comme le montre la complétion du carré :
\[
 \left\langle\binom b m,Y\binom b m\right\rangle
 =\langle b,Sb\rangle+
 \langle m+D_i^{-1}C^*b,D_i(m+D_i^{-1}C^*b)\rangle.
\]
Par conséquent, la réduction cohérente donne
\[
 H_{\rm eff}=\frac{\hbar c}{L_\alpha}S,\qquad
 R_{\rm eff}=L_\alpha S,\qquad\Lambda_b=L_\alpha S^{-1}.
\]
La compression de $Y^{-1}$ sur le bloc de lecture est $S^{-1}$ ;
la compression directe de $Y$ est $A$. La mémoire distingue ces opérations.

Si l’élimination dynamique produit une métrique temporelle bornée
$Z\ge z_0I$, $z_0>0$, on prend $T=Z^{1/2}$, borné et inversible.
La normalisation canonique exige les transports duaux
\[
 H_c=T^{-*}H_{\rm eff}T^{-1},\qquad
 R_c=T^{-*}R_{\rm eff}T^{-1},\qquad
 \Lambda_c=T\Lambda_bT^*.
\]
Leur multiplication donne $H_c\Lambda_c=\hbar cI$ et
$R_c\Lambda_c=L_\alpha^2I$, sans hypothèse de commutativité.
Dans le développement résolvant d’une énergie de blocs $H_{bb},V,H_{ii}$,
avec $H_{ii}$ positif inversible et
$|\hbar\omega|<\|H_{ii}^{-1}\|^{-1}$,
\[
 K_b(\omega)=H_{bb}-\hbar\omega I
 -V(H_{ii}-\hbar\omega I)^{-1}V^*
 =H_{\rm stat}-\hbar\omega Z+O(\omega^2),
 \qquad Z=I+VH_{ii}^{-2}V^*.
\]
La dernière expression est locale autour de zéro et conserve l’ordre
d’approximation ; les produits duaux précédents sont des identités exactes
des lecteurs normalisés.

La proportionnalité est également conservée dans l’élimination dynamique
complète. Pour $z$ tel que $H_{ii}-zI$ soit inversible, nous définissons
\[
 \mathcal K_H(z)=H_{bb}-zI-V(H_{ii}-zI)^{-1}V^*,\qquad
 \kappa=\frac{L_\alpha^2}{\hbar c}.
\]
La même décomposition de lecture et de mémoire appliquée à $R=\kappa H$
produit
\begin{equation}
 \mathcal K_R(\kappa z)=\kappa\mathcal K_H(z).
 \label{md:rigidez:resolvente-exacto}
\end{equation}
En effet, ses blocs sont $R_{bb}=\kappa H_{bb}$,
$R_{bi}=\kappa V$ et $R_{ii}=\kappa H_{ii}$, et
\[
 (R_{ii}-\kappa zI)^{-1}
 =\kappa^{-1}(H_{ii}-zI)^{-1}.
\]
La substitution dans le complément de Schur démontre \eqref{md:rigidez:resolvente-exacto}, en conservant l'ordre des facteurs. Lorsque les deux compléments sont inversibles, leurs inverses satisfont $\mathcal K_R(\kappa z)^{-1}=\kappa^{-1}\mathcal K_H(z)^{-1}$. Le paramètre énergétique est $z=\hbar\omega$ et son image dans le lecteur radial est $\kappa z$ : tous deux conservent leurs dimensions respectives. L'identité vaut sur tout ce domaine résolvant et transporte donc conjointement tous les ordres de la mémoire dynamique. Le développement local précédent décrit ses premiers coefficients ; cette égalité explique la conservation exacte du rapport radial lors de la réduction.

\subsection{Rigidité temporelle et mémoire de frontière}
\label{md:rigidez:memoria}

La conservation de la seule aire admet $R\mapsto\zeta R$, $\Lambda\mapsto\Lambda/\zeta$. Conserver également $H,\hbar,c$ et $H\Lambda=\hbar cI$ impose $\zeta=1$. Si $H\mapsto\zeta H$ est également transformé, le quotient $R/H$ et le coefficient $G$ restent égaux. Les comparaisons doivent conserver les mêmes observations conjointes.

La phase relevée apporte une autre restriction effective. Dans une réalisation bornée $H=H_{\rm clk}+D_0^*WD_0$ avec $W>0$, l'évolution $P(t)=\exp(-itH/\hbar)$ détermine $H=i\hbar\dot P(0)$. Une fois $H_{\rm clk},D_0,\hbar$ et l'horloge fixés, l'égalité des évolutions implique $D_0^*(W'-W)D_0=0$. Pour un parcours d'extrémité initiale fixée, $D_0$ est triangulaire inversible et l'on en déduit $W'=W$. Pour un graphe général, on détermine la forme sur $\operatorname{Ran}D_0$ ; étendre ensuite le domaine des résidus modifierait le problème d'élimination. Un seul retour $P(\tau)$ admet des alias spectraux, tandis que la phase relevée ou les temps compatibles tendant vers zéro les distinguent.

Dans la coordinatisation de capacité nonadique, soient $\rho=\log_{10}9$, $\delta=1-\rho$, $p_j=\lfloor j/\delta\rfloor$ et $v_j=p_j-j$. La capacité d'une fenêtre est
\[
 C(j,n)=p_{j+n}-p_j-n.
\]
Pour $n=35$, ses valeurs sont $729$ ou $730$ selon l'origine ; le retour de capacité modulo neuf distingue la première au sein de cette classe. La fenêtre $j=20\to55$ satisfait $p_j:437\to1201$ et $v_j:417\to1146$. Sa phase de capacité revient à $3$ et le quotient augmente de $46$ à $127$. Les six capacités de $20$ et les vingt-neuf de $21$ ont pour somme $729$ ; la mémoire accumulée fait toujours partie de l'état.

Pour la première arrivée $T(k)=\lceil k/\rho\rceil-1$, le dénombrement est $N(k,C)=T(k+C)-T(k)-C$. Écrire $\epsilon(k)=\lceil k/\rho\rceil-k/\rho$ donne exactement
\[
 N(k,C)=(\rho^{-1}-1)C+\epsilon(k+C)-\epsilon(k).
\]
Les fenêtres de capacité $729$ dont les dénombrements sont $34$ et $35$ conservent des termes de frontière différents. Même $k=21$ et $k=2289$ partagent $(k\bmod108,T(k)\bmod108)=(21,22)$ et accumulent respectivement $1$ et $109$ vacances. La coïncidence des calendriers finis conserve moins d'information que l'état complet. Ces calculs interprètent le retour ; le lecteur longitudinal de \eqref{md:rigidez:longitud-def} possède son propre constructeur et ne sélectionne pas son échelle à l'aide de ces fenêtres.

\subsection{Normalisation géométrique et évaluation}
\label{md:rigidez:evaluacion}

Le bloc central d'APP
$B_c=\left(\begin{smallmatrix}7&2\\2&7\end{smallmatrix}\right)$
satisfait $B_c^{\mathsf T}g_cB_c=45g_c$, $g_c=\operatorname{diag}(1,-1)$. Sa normalisation positive unimodulaire est $B_c/\sqrt{45}$ : un facteur positif supplémentaire $a$ donnerait le déterminant $a^2$. L'incidence aréale--volumétrique
$T_{\rm av}=\left(\begin{smallmatrix}1&1\\1&2\end{smallmatrix}\right)$
conserve
$G_{\rm av}=\left(\begin{smallmatrix}2&1\\1&-2\end{smallmatrix}\right)$
par multiplication directe. Dans une réalisation de Minkowski, la forme causale et la soudure dimensionnelle jouent des rôles distincts. La cochaîne $\beta_m=\ell_m\sigma_mW_\square e_\lambda$, $\ell_m=\ell_0/9^m$, et la relation $\ell_0=ct_0$ conservent l'échelle avec le cône. Modifier la soudure en ne conservant que la signature change cette réalisation.

L'unité expansive de Pell de norme un et de trace quatre satisfait $x+x^{-1}=4$, $x>1$, d'où $x=2+\sqrt3$. L'orbite $r_n=r_0x^n$ conserve $r_nr_{-n}=r_0^2$. Lorsque $r_0=L_\alpha$, le produit conserve la même section aréale. La norme, la forme marquée, le retour et la section longitudinale restreignent des objets différents ; leurs fonctions sont maintenues lors de leur composition.

\Needspace{10\baselineskip}
Dans la section SI déclarée, $L_*=1\,\mathrm m$, $S_*=10^{-34}\,\mathrm{J\,s}$ et $c=299792458\,\mathrm{m\,s^{-1}}$, la substitution structurelle produit
\[
 L_\alpha=1.616254982625126330\ldots\,10^{-35}\,\mathrm m,\qquad
 \hbar_{\rm ret}=1.054571816915039061\ldots\,10^{-34}\,\mathrm{J\,s},
\]
\begin{equation}
 G_{\rm ret}=6.674299659414022706\ldots\,10^{-11}
 \,\mathrm{m^3\,kg^{-1}\,s^{-2}}.
 \label{md:rigidez:valor}
\end{equation}
L'évaluateur substitue les coordonnées régionales et le registre $K$, reconstruit $\alpha$, le retour de $H_5$ et $L_\alpha$ ; la comparaison métrologique intervient ensuite. La référence CODATA 2022 est $6.67430(15)\,10^{-11}$ dans ces unités \cite{codata2022}. La différence est $-3.40586\,10^{-18}$, soit $-0.00227$ incertitudes-types. Les chiffres calculés décrivent l'évaluation mathématique dans cette section, et non une incertitude expérimentale. Les sections d'action antérieure et retournée conservent leurs étiquettes et $\hbar_aG_a=c^3L_\alpha^2$.

La signification structurelle de $G$ est ainsi déterminée : c'est le coefficient qui convertit, dans la réalisation radiale commune, la lecture énergétique en longueur conjuguée en préservant simultanément l'action et l'aire. La mémoire, les marques de feuillet et les bases font partie de la construction qui fixe ce coefficient. Leur conservation conjointe explique l'unicité démontrée.
